% !PRINTED 328
\weberpage{328}
%<<CONT
genügt. Wenn insbesondere die Form von den Variablen $x$ unabhängig und nur von den Coëfficienten $a$ abhängig ist, so wird sie zur \gesperrt{Invariante}. Der Exponent $\lambda$, der das \gesperrt{Gewicht} der Covariante heisst, ist eine ganze Zahl.

Wir haben im Bd. I, \S.\,59 eine bei allen Formen, deren Grad grösser als 1 ist, vorhandene Covariante vom Gewicht 2 kennen gelernt, nämlich die \gesperrt{Hesse'sche Determinante}:

\weq{2}{H(x_1, x_2, x_3) = \begin{vmatrix} f''(x_1, x_1), & f''(x_1, x_2), & f''(x_1, x_3) \\ f''(x_2, x_1), & f''(x_2, x_2), & f''(x_2, x_3) \\ f''(x_3, x_1), & f''(x_3, x_2), & f''(x_3, x_3) \end{vmatrix}.}

Ausser dieser Covariante $H$ wollen wir noch zwei andere allgemeine Covarianten betrachten, die bei allen ternären Formen von höherem als dem $2^{\text{ten}}$ Grade existiren, und die sich leicht auch für Formen von beliebig vielen Variablen verallgemeinern lassen. Die erste dieser Formen ist, wenn $H$ wie oben die \gesperrt{Hesse}'sche Covariante bedeutet:

\weq{3}{C_1(x, a) = \begin{vmatrix} f''(x_1, x_1), & f''(x_1, x_2), & f''(x_1, x_3), & H'(x_1) \\ f''(x_2, x_1), & f''(x_2, x_2), & f''(x_2, x_3), & H'(x_2) \\ f''(x_3, x_1), & f''(x_3, x_2), & f''(x_3, x_3), & H'(x_3) \\ H'(x_1), & H'(x_2), & H'(x_3) & 0 \end{vmatrix}.}

Um die Invarianten-Eigenschaft für diese Form nachzuweisen, machen wir darin die Substitution \S.\,85, (1). Dadurch geht $f(x, a)$ in $\varphi(y, b)$ und $r^2 H(x, a)$ in $H(y, b)$ über. Wenn wir dann unter $H'(y_1), H'(y_2), H'(y_3)$ die Derivirten von $H(y, b)$ verstehen, so haben wir die Formeln

\weq{4}{
\begin{aligned}
\varphi''(y_\mu, y_\nu) &= \sum\limits^{i,k}_{i} f''(x_i, x_k)\, \alpha_i^{(\mu)} \alpha_k^{(\nu)} \\
H'(y_\nu) &= r^2 \sum H'(x_i)\, \alpha_i^{(\nu)},
\end{aligned}
}
in denen $\mu, \nu, i, k$ von 1 bis 3 laufen.

Multiplicirt man nun die Determinante $C_1(x, a)$ zweimal nach einander mit der in der Form

\[
r = \begin{vmatrix} \alpha_1^{(1)}, & \alpha_1^{(2)}, & \alpha_1^{(3)}, & 0 \\ \alpha_2^{(1)}, & \alpha_2^{(2)}, & \alpha_2^{(3)}, & 0 \\ \alpha_3^{(1)}, & \alpha_3^{(2)}, & \alpha_3^{(3)}, & 0 \\ 0, & 0, & 0, & 1 \end{vmatrix}
\]
geschriebenen Substitutionsdeterminante, indem man das eine
%CONT>>
